\subsection{Positive Semidefinite and Positive Definite Matrices}

\begin{definition} \label{def:positive-semidefinite-matrix} \hlt{Positive Semidefinite and Positive Definite Matrices}\\
Let $n\in\mathbb{Z}_{>0}$, $\mathbb{F}\in\{\mathbb{R},\mathbb{C}\}$, and $A\in\operatorname{Mat}(n,n,\mathbb{F})$.
\begin{enumerate}[label=(\roman*)]
\setlength{\itemsep}{0pt}
\item The matrix $A$ is \hlt{positive semidefinite}, written $A\succeq0$, iff $A=A^*$ and $(\forall\,\mathbf{x}\in\mathbb{F}^n)\,\mathbf{x}^*A\mathbf{x}\geq0$.
\item The matrix $A$ is \hlt{positive definite}, written $A\succ0$, iff $A=A^*$ and $(\forall\,\mathbf{x}\in\mathbb{F}^n\setminus\{\mathbf{0}\})\,\mathbf{x}^*A\mathbf{x}>0$.
\item The matrix $A$ is \hlt{negative semidefinite} iff $-A\succeq0$.
\item The matrix $A$ is \hlt{negative definite} iff $-A\succ0$.
\end{enumerate}
\end{definition}

\begin{proposition} \label{prop:rectangular-gram-products-are-psd} \hlt{Rectangular Gram Products Are Positive Semidefinite}\\
Let $m,n\in\mathbb{Z}_{>0}$, $\mathbb{F}\in\{\mathbb{R},\mathbb{C}\}$, and $R\in\operatorname{Mat}(m,n,\mathbb{F})$.
\begin{enumerate}[label=(\roman*)]
\setlength{\itemsep}{0pt}
\item $R^*R\succeq0$.
\item $RR^*\succeq0$.
\item $(\forall\,j\in\{1,\ldots,n\})\,(R^*R)_{j,j}=\norm{R_{\cdot,j}}_2^2$.
\item $(\forall\,i\in\{1,\ldots,m\})\,(RR^*)_{i,i}=\norm{R_{i,\cdot}}_2^2$.
\item If $\operatorname{rank}(R)=n$, then $R^*R\succ0$.
\item If $\operatorname{rank}(R)=m$, then $RR^*\succ0$.
\end{enumerate}
\begin{proof}
For $(A,d)\in\{(R,n),(R^*,m)\}$, $(A^*A)^*=A^*A$ by Proposition \ref{prop:basic-adjoint-identities}(ii),(iii).\\
For the same pairs, $(\forall\,\mathbf{z}\in\mathbb{F}^d)\,\mathbf{z}^*A^*A\mathbf{z}=\norm{A\mathbf{z}}_2^2\geq0$.
\begin{enumerate}[label=(\roman*)]
\setlength{\itemsep}{0pt}
\item $R^*R\succeq0$ by the preceding identities with $(A,d)=(R,n)$.
\item $RR^*\succeq0$ by the preceding identities with $(A,d)=(R^*,m)$.
\item $(\forall\,j\in\{1,\ldots,n\})\,(R^*R)_{j,j}=\mathbf{e}_j^*R^*R\mathbf{e}_j=\norm{R\mathbf{e}_j}_2^2=\norm{R_{\cdot,j}}_2^2$.
\item $(\forall\,i\in\{1,\ldots,m\})\,(RR^*)_{i,i}=\mathbf{e}_i^*RR^*\mathbf{e}_i=\norm{R^*\mathbf{e}_i}_2^2=\norm{R_{i,\cdot}}_2^2$.
\item If $\operatorname{rank}(R)=n$, then $\ker(R)=\{\mathbf{0}\}$ by Corollary \ref{cor:matrix-rank-nullity}(i). Thus $(\forall\,\mathbf{x}\in\mathbb{F}^n\setminus\{\mathbf{0}\})\,\mathbf{x}^*R^*R\mathbf{x}=\norm{R\mathbf{x}}_2^2>0$, and $R^*R\succ0$ by part (i).
\item If $\operatorname{rank}(R)=m$, then $\operatorname{col}(R)=\mathbb{F}^m$ by rank definition, Proposition \ref{prop:dimension-bounds-and-basis-criteria}(ii); $\ker(R^*)=\operatorname{col}(R)^\perp=\{\mathbf{0}\}$ by Proposition \ref{prop:adjoint-kernel-range-orthogonality}(i).\\
Therefore $(\forall\,\mathbf{y}\in\mathbb{F}^m\setminus\{\mathbf{0}\})\,\mathbf{y}^*RR^*\mathbf{y}=\norm{R^*\mathbf{y}}_2^2>0$, and $RR^*\succ0$ by part (ii).
\end{enumerate}
\hlt{Structural Commentary}: Adjoint identities establish self-adjointness; rank removes the kernels in the two squared-norm forms.\\
\hlt{Key Insight}: A Gram form is strictly positive exactly off its factor's kernel.
\end{proof}
\end{proposition}

\begin{theorem} \label{thm:positive-semidefinite-matrix-equivalences} \hlt{Equivalent Forms of Positive Semidefiniteness}\\
Let $n\in\mathbb{Z}_{>0}$, $\mathbb{F}\in\{\mathbb{R},\mathbb{C}\}$, and $A\in\operatorname{Mat}(n,n,\mathbb{F})$. The following are equivalent:
\begin{enumerate}[label=(\roman*)]
\setlength{\itemsep}{0pt}
\item $A\succeq0$.
\item $A=A^*$ and $\operatorname{spec}(A)\subseteq[0,\infty)$.
\item $(\exists\,C\in\operatorname{Mat}(n,n,\mathbb{F}))\,[C\succeq0\land C^2=A]$.
\item $(\exists\,m\in\mathbb{Z}_{>0})(\exists\,B\in\operatorname{Mat}(m,n,\mathbb{F}))\,A=B^*B$.
\end{enumerate}
\begin{proof}
Identify each square matrix with its induced operator on $\mathbb{F}^n$.
\begin{enumerate}[label=(\roman*)]
\setlength{\itemsep}{0pt}
\item[$(i)$]$\Rightarrow (ii)$ $A=A^*$ and $\operatorname{spec}(A)\subseteq[0,\infty)$ by Theorem \ref{thm:positive-operator-characterization}(i)$\Rightarrow$(ii). \item[$(ii)$]$\Rightarrow (iii)$ Choose $C\succeq0$ with $C^2=A$ by Theorem \ref{thm:positive-operator-characterization}(ii)$\Rightarrow$(iii). \item[$(iii)$]$\Rightarrow (iv)$ If $C\succeq0$ and $C^2=A$, then $A=C^2=C^*C$; take $(m,B)=(n,C)$. \item[$(iv)$]$\Rightarrow (i)$ $A\succeq0$ by Proposition \ref{prop:rectangular-gram-products-are-psd}(i) for $A=B^*B$.
\end{enumerate}
\hlt{Structural Commentary}: Operator positivity supplies the spectral root; self-adjointness converts that root into a Gram factor.\\
\hlt{Key Insight}: For $C\succeq0$, the identities $C^2=C^*C$ coincide.
\end{proof}
\end{theorem}

\begin{corollary} \label{cor:positive-semidefinite-square-root-formula} \hlt{Unique Positive Semidefinite Square Root}\\
Let $n\in\mathbb{Z}_{>0}$, $\mathbb{F}\in\{\mathbb{R},\mathbb{C}\}$, and $A\in\operatorname{Mat}(n,n,\mathbb{F})$ satisfy $A\succeq0$.
\begin{enumerate}[label=(\roman*)]
\setlength{\itemsep}{0pt}
\item $(\exists!\,C\in\operatorname{Mat}(n,n,\mathbb{F}))\,[C\succeq0\land C^2=A]$; this matrix is denoted by $A^{1/2}$.
\item If $A=U\operatorname{diag}(\lambda_1,\ldots,\lambda_n)U^*$, where $U$ is unitary and $(\forall\,j\in\{1,\ldots,n\})\,\lambda_j\geq0$,\\
then $A^{1/2}=U\operatorname{diag}(\sqrt{\lambda_1},\ldots,\sqrt{\lambda_n})U^*$.
\end{enumerate}
\begin{proof}
\begin{enumerate}[label=(\roman*)]
\setlength{\itemsep}{0pt}
\item Existence and uniqueness hold by Proposition \ref{prop:positive-square-root-unique} for the operator induced by $A$.
\item Set $C\defeq U\operatorname{diag}(\sqrt{\lambda_1},\ldots,\sqrt{\lambda_n})U^*$. Then $C=C^*$, $\operatorname{spec}(C)\subseteq[0,\infty)$, and $C^2=A$.\\
Thus $C\succeq0$ by Theorem \ref{thm:positive-semidefinite-matrix-equivalences}(ii)$\Rightarrow$(i), and $C=A^{1/2}$ by part (i).
\end{enumerate}
\hlt{Structural Commentary}: Spectral coordinates construct the root; operator uniqueness identifies every positive semidefinite root.\\
\hlt{Key Insight}: Each nonnegative eigenvalue has one nonnegative square root.
\end{proof}
\end{corollary}

\begin{theorem} \label{thm:gram-factor-kernel-rank-column-data} \hlt{Kernel, Rank, and Columns of $B^*B$}\\
Let $m,n\in\mathbb{Z}_{>0}$, $\mathbb{F}\in\{\mathbb{R},\mathbb{C}\}$, $B\in\operatorname{Mat}(m,n,\mathbb{F})$, and $A=B^*B$.
\begin{enumerate}[label=(\roman*)]
\setlength{\itemsep}{0pt}
\item $A\succeq0$;
\item $\ker(A)=\ker(B)$;
\item $\operatorname{rank}(A)=\operatorname{rank}(B)$;
\item $\operatorname{col}(A)=\operatorname{col}(B^*)$;
\item $A=0$ iff $B=0$.
\end{enumerate}
\begin{proof}
\begin{enumerate}[label=(\roman*)]
\setlength{\itemsep}{0pt}
\item[$(i)$] $A\succeq0$ by Proposition \ref{prop:rectangular-gram-products-are-psd}(i).
\item[$(ii)$] $\ker(B)\subseteq\ker(A)$ by $A=B^*B$. If $\mathbf{x}\in\ker(A)$, then $\norm{B\mathbf{x}}_2^2=\mathbf{x}^*A\mathbf{x}=0$, so $\mathbf{x}\in\ker(B)$.
\item[$(iii)$] $\operatorname{rank}(A)=n-\dim\ker(A)=n-\dim\ker(B)=\operatorname{rank}(B)$ by part (ii) and Corollary \ref{cor:matrix-rank-nullity}(i).
\item[$(iv)$] $\operatorname{col}(A)=\ker(A)^\perp=\ker(B)^\perp=\operatorname{col}(B^*)$ by parts (i)--(ii) and Proposition \ref{prop:adjoint-kernel-range-orthogonality}(iii).
\item[$(v)$] ($\Rightarrow$) If $A=0$, then $\operatorname{rank}(B)=\operatorname{rank}(A)=0$ by part (iii), so $B=0$.
\item[$(v)$] ($\Leftarrow$) $B=0\Rightarrow A=B^*B=0$.
\end{enumerate}
\hlt{Structural Commentary}: Squared-norm identity fixes the kernel; rank--nullity and adjoint orthogonality determine rank and columns.\\
\hlt{Key Insight}: $B^*B$ and $B$ vanish on exactly the same vectors.
\end{proof}
\end{theorem}

Figure \ref{fig:psd-factorization-kernel} summarizes the squared-norm argument that transfers kernel data from $B^*B$ to $B$.

\begin{illustrationfigure}[H]
  \centering
  \begin{tikzpicture}[
  >=Stealth,
  every node/.style={font=\small},
  box/.style={draw=gray!55, rounded corners=3pt, fill=gray!6, align=center, inner xsep=6pt, inner ysep=4.5pt, minimum height=0.82cm},
  result/.style={box, fill=green!7, minimum width=2.4cm},
  arrow/.style={->, line width=0.85pt, blue!70!black}
]
  \node[box, minimum width=2.05cm] (factor) at (0,0) {$A=B^*B$};
  \node[box, minimum width=4.0cm] (test) at (3.55,0)
    {$\mathbf{x}^*A\mathbf{x}=\norm{B\mathbf{x}}_2^2$\\[-2pt]
     $\norm{B\mathbf{x}}_2^2=0\Longleftrightarrow B\mathbf{x}=0$};
  \node[result] (kernel) at (7.2,0) {$\ker(A)=\ker(B)$};

  \draw[arrow] (factor) -- (test);
  \draw[arrow] (test) -- (kernel);
\end{tikzpicture}

  \caption{For $A=B^*B$, the squared-norm identity certifies $\ker(A)=\ker(B)$.}
  \label{fig:psd-factorization-kernel}
\end{illustrationfigure}

\begin{corollary} \label{cor:positive-semidefinite-zero-quadratic-form} \hlt{Zero Quadratic Form Detects the Kernel}\\
Let $n\in\mathbb{Z}_{>0}$, $\mathbb{F}\in\{\mathbb{R},\mathbb{C}\}$, and $A\in\operatorname{Mat}(n,n,\mathbb{F})$ satisfy $A\succeq0$. Then $(\forall\,\mathbf{x}\in\mathbb{F}^n)\,[\mathbf{x}^*A\mathbf{x}=0\Longleftrightarrow A\mathbf{x}=\mathbf{0}]$.
\begin{proof}
Fix $\mathbf{x}\in\mathbb{F}^n$, and choose $m\in\mathbb{Z}_{>0}$ and $B\in\operatorname{Mat}(m,n,\mathbb{F})$ with $A=B^*B$ by Theorem \ref{thm:positive-semidefinite-matrix-equivalences}(iv).
\begin{enumerate}[label=\roman*.]
\setlength{\itemsep}{0pt}
\item[$(\Rightarrow)$] $\mathbf{x}^*A\mathbf{x}=0\Rightarrow\norm{B\mathbf{x}}_2^2=0\Rightarrow B\mathbf{x}=\mathbf{0}\Rightarrow A\mathbf{x}=\mathbf{0}$.
\item[$(\Leftarrow)$] $A\mathbf{x}=\mathbf{0}\Rightarrow\mathbf{x}^*A\mathbf{x}=0$.
\end{enumerate}
\hlt{Structural Commentary}: Gram factorization converts the quadratic value into a norm whose zero set is $\ker(A)$.\\
\hlt{Key Insight}: $\norm{B\mathbf{x}}_2^2=0$ forces $B\mathbf{x}=\mathbf{0}$.
\end{proof}
\end{corollary}

\begin{computationmethod} \label{meth:psd-spectral-factorization} \hlt{PSD Spectral Factorization}\\
Let $n\in\mathbb{Z}_{>0}$, $\mathbb{F}\in\{\mathbb{R},\mathbb{C}\}$, and let $A\in\operatorname{Mat}(n,n,\mathbb{F})$ satisfy $A=A^*$.\\
The procedure decides whether $A\succeq0$, constructs $A^{1/2}$ when it exists, and transfers kernel and rank questions to that factor.
\begin{enumerate}[label=(\roman*)]
\setlength{\itemsep}{0pt}
\item Compute a unitary diagonalization $A=U\operatorname{diag}(\lambda_1,\ldots,\lambda_n)U^*$ by the appropriate spectral theorem.
\item If $(\exists\,j\in\{1,\ldots,n\})\,\lambda_j<0$, conclude that $A\not\succeq0$. If $(\forall\,j\in\{1,\ldots,n\})\,\lambda_j\geq0$, conclude that $A\succeq0$ by Theorem \ref{thm:positive-semidefinite-matrix-equivalences}(ii).
\item In the positive semidefinite case, form $C=U\operatorname{diag}(\sqrt{\lambda_1},\ldots,\sqrt{\lambda_n})U^*$. Then $C=A^{1/2}$ and $A=C^*C$ by Corollary \ref{cor:positive-semidefinite-square-root-formula}(i),(ii).
\item Read the kernel, rank, column-space, and zero-quadratic-form data from $C$ by Theorem \ref{thm:gram-factor-kernel-rank-column-data} and Corollary \ref{cor:positive-semidefinite-zero-quadratic-form}.
\end{enumerate}
\end{computationmethod}

\begin{theorem} \label{thm:positive-definite-matrix-equivalences} \hlt{Equivalent Forms of Positive Definiteness}\\
Let $n\in\mathbb{Z}_{>0}$, $\mathbb{F}\in\{\mathbb{R},\mathbb{C}\}$, and $A\in\operatorname{Mat}(n,n,\mathbb{F})$. The following are equivalent:
\begin{enumerate}[label=(\roman*)]
\setlength{\itemsep}{0pt}
\item $A\succ0$.
\item $A=A^*$ and $\operatorname{spec}(A)\subseteq(0,\infty)$.
\item $A\succeq0$ and $A$ is invertible.
\item $(\exists\,B\in\operatorname{Mat}(n,n,\mathbb{F})^\times)\,A=B^*B$.
\item $A$ is invertible and $A^{-1}\succ0$.
\end{enumerate}
\begin{proof}
\begin{enumerate}[label=(\roman*)]
\setlength{\itemsep}{0pt}
\item[$(i)$]$\Rightarrow (ii)$ $A=A^*$ by positive definiteness. For $\lambda\in\operatorname{spec}(A)$ and $\mathbf{v}\in E_\lambda(A)\setminus\{\mathbf{0}\}$, $0<\mathbf{v}^*A\mathbf{v}=\lambda\norm{\mathbf{v}}_2^2$, so $\lambda>0$. \item[$(ii)$]$\Rightarrow (iii)$ $A\succeq0$ by Theorem \ref{thm:positive-semidefinite-matrix-equivalences}(ii)$\Rightarrow$(i), and $A$ is invertible by Corollary \ref{cor:matrix-invertible-iff-zero-not-spectrum}. 
\item[$(iii)$]$\Rightarrow (iv)$ Set $C\defeq A^{1/2}$. Then $A=C^2=C^*C$ by Corollary \ref{cor:positive-semidefinite-square-root-formula}(i), while $\ker(C)\subseteq\ker(A)=\{\mathbf{0}\}$.\\
Thus $C$ is invertible by Corollary \ref{cor:square-matrix-algebraic-invertibility-dictionary}(vi)$\Rightarrow$(i). 
\item[$(iv)$]$\Rightarrow (v)$ $A=B^*B$ is invertible and $A^{-1}=B^{-1}(B^{-1})^*$ by Corollary \ref{cor:matrix-inverse-product-and-transpose}(i) and Proposition \ref{prop:adjoints-inverses-and-polynomials}(i).\\
The equality $\operatorname{rank}(B^{-1})=n$ holds by Corollary \ref{cor:square-matrix-invertible-iff-full-rank}, so $A^{-1}\succ0$ by Proposition \ref{prop:rectangular-gram-products-are-psd}(vi) with $R=B^{-1}$.
 \item[$(v)$]$\Rightarrow (i)$ $(A^{-1})^*=A^{-1}$ by positive definiteness, so $A^*=((A^{-1})^*)^{-1}=A$ by Proposition \ref{prop:adjoints-inverses-and-polynomials}(i).\\
For $\mathbf{x}\in\mathbb{F}^n\setminus\{\mathbf{0}\}$, set $\mathbf{y}\defeq A\mathbf{x}\neq\mathbf{0}$. Then $0<\mathbf{y}^*A^{-1}\mathbf{y}=\mathbf{x}^*A\mathbf{x}$, so $A\succ0$.
\end{enumerate}
\hlt{Structural Commentary}: Spectral positivity removes zero, the square root supplies a Gram factor, and substitution reverses inversion.\\
\hlt{Key Insight}: $\mathbf{y}=A\mathbf{x}$ transfers strict positivity from $A^{-1}$ to $A$.
\end{proof}
\end{theorem}

\begin{proposition} \label{prop:trace-determinant-positive-semidefinite} \hlt{Trace and Determinant of Positive Semidefinite Matrices}\\
Let $n\in\mathbb{Z}_{>0}$, $\mathbb{F}\in\{\mathbb{R},\mathbb{C}\}$, and $A\in\operatorname{Mat}(n,n,\mathbb{F})$ satisfy $A\succeq0$.
\begin{enumerate}[label=(\roman*)]
\setlength{\itemsep}{0pt}
\item $\operatorname{tr}(A)\geq0$.
\item $\operatorname{tr}(A)=0$ iff $A=0$.
\item $\det(A)\geq0$.
\item $\det(A)>0$ iff $A\succ0$.
\end{enumerate}
\begin{proof}
Write $A=U\operatorname{diag}(\lambda_1,\ldots,\lambda_n)U^*$ by the appropriate one of Theorems \ref{thm:complex-spectral-theorem} and \ref{thm:real-spectral-theorem}.\\
The relation $(\forall\,j\in\{1,\ldots,n\})\,\lambda_j\geq0$ holds by Theorem \ref{thm:positive-semidefinite-matrix-equivalences}(ii). Then $\operatorname{tr}(A)=\sum\limits_{j=1}^{n}\lambda_j$ and $\det(A)=\prod\limits_{j=1}^{n}\lambda_j$ by Theorem \ref{thm:similarity-preserves-trace-and-determinant}(i),(ii).
\begin{enumerate}[label=(\roman*)]
\setlength{\itemsep}{0pt}
\item[$(i)$] $\operatorname{tr}(A)=\sum\limits_{j=1}^{n}\lambda_j\geq0$.
\item[$(ii)$] ($\Rightarrow$) $\operatorname{tr}(A)=0\Rightarrow(\forall\,j\in\{1,\ldots,n\})\,\lambda_j=0\Rightarrow A=0$.
\item[$(ii)$] ($\Leftarrow$) $A=0\Rightarrow\operatorname{tr}(A)=0$.
\item[$(iii)$] $\det(A)=\prod\limits_{j=1}^{n}\lambda_j\geq0$.
\item[$(iv)$] ($\Rightarrow$) $\det(A)>0\Rightarrow(\forall\,j\in\{1,\ldots,n\})\,\lambda_j>0$, so $A\succ0$ by Theorem \ref{thm:positive-definite-matrix-equivalences}(ii).
\item[$(iv)$] ($\Leftarrow$) $A\succ0\Rightarrow(\forall\,j\in\{1,\ldots,n\})\,\lambda_j>0$ by Theorem \ref{thm:positive-definite-matrix-equivalences}(ii), so $\det(A)>0$.
\end{enumerate}
\hlt{Structural Commentary}: Unitary diagonalization transfers trace and determinant to a nonnegative eigenvalue sum and product.\\
\hlt{Key Insight}: The sum detects universal vanishing; the product detects any vanishing.
\end{proof}
\end{proposition}

\begin{theorem} \label{thm:positive-semidefinite-stability-rules} \hlt{Stability Rules for Positive Semidefinite Matrices}\\
Let $m,n,p\in\mathbb{Z}_{>0}$, $\mathbb{F}\in\{\mathbb{R},\mathbb{C}\}$. Let $A,C\in\operatorname{Mat}(n,n,\mathbb{F})$, $B\in\operatorname{Mat}(m,m,\mathbb{F})$,  $S\in\operatorname{Mat}(n,p,\mathbb{F})$. Fix $a,b\in\mathbb{R}_{\geq0}$, suppose $A,C,B\succeq0$.
\begin{enumerate}[label=(\roman*)]
\setlength{\itemsep}{0pt}
\item $aA+bC\succeq0$.
\item $S^*AS\succeq0$.
\item $A\oplus B\succeq0$.
\item If $a,b>0$ and either $A\succ0$ or $C\succ0$, then $aA+bC\succ0$.
\item $S^*AS\succ0$ iff $(\forall\,\mathbf{x}\in\mathbb{F}^p\setminus\{\mathbf{0}\})\,S\mathbf{x}\notin\ker(A)$.
\item $A\oplus B\succ0$ iff $A\succ0$ and $B\succ0$.
\end{enumerate}
\begin{proof}
For a square matrix $M$ and compatible vector $\mathbf{z}$, set $q_M(\mathbf{z})\defeq\mathbf{z}^*M\mathbf{z}$.
\begin{enumerate}[label=(\roman*)]
\setlength{\itemsep}{0pt}
\item[$(i)$] $(aA+bC)^*=aA+bC$ and $(\forall\,\mathbf{x}\in\mathbb{F}^n)\,q_{aA+bC}(\mathbf{x})=aq_A(\mathbf{x})+bq_C(\mathbf{x})\geq0$, so $aA+bC\succeq0$.
\item[$(ii)$] $(S^*AS)^*=S^*AS$ by Proposition \ref{prop:basic-adjoint-identities}(ii),(iii). The identity $(\forall\,\mathbf{x}\in\mathbb{F}^p)\,q_{S^*AS}(\mathbf{x})=q_A(S\mathbf{x})\geq0$ implies $S^*AS\succeq0$.
\item[$(iii)$] $(A\oplus B)^*=A\oplus B$, and $(\forall\,(\mathbf{x},\mathbf{y})\in\mathbb{F}^n\times\mathbb{F}^m)\,q_{A\oplus B}(\mathbf{x}\oplus\mathbf{y})=q_A(\mathbf{x})+q_B(\mathbf{y})$. The right side is nonnegative, so $A\oplus B\succeq0$.
\item[$(iv)$] If $a,b>0$ and $A\succ0$, then $(\forall\,\mathbf{x}\in\mathbb{F}^n\setminus\{\mathbf{0}\})\,aq_A(\mathbf{x})>0$ and $bq_C(\mathbf{x})\geq0$, so $aA+bC\succ0$ by part (i).\\
If $a,b>0$ and $C\succ0$, then $(\forall\,\mathbf{x}\in\mathbb{F}^n\setminus\{\mathbf{0}\})\,aq_A(\mathbf{x})\geq0$ and $bq_C(\mathbf{x})>0$, so again $aA+bC\succ0$.
\item[$(v)$] ($\Rightarrow$) If $S^*AS\succ0$ and $\mathbf{x}\in\mathbb{F}^p\setminus\{\mathbf{0}\}$, then $q_A(S\mathbf{x})=q_{S^*AS}(\mathbf{x})>0$. Thus $S\mathbf{x}\notin\ker(A)$ by Corollary \ref{cor:positive-semidefinite-zero-quadratic-form}.
\item[$(v)$] ($\Leftarrow$) If $\mathbf{x}\in\mathbb{F}^p\setminus\{\mathbf{0}\}$, then $S\mathbf{x}\notin\ker(A)$ by hypothesis.\\
Thus $q_{S^*AS}(\mathbf{x})=q_A(S\mathbf{x})>0$ by $A\succeq0$ and Corollary \ref{cor:positive-semidefinite-zero-quadratic-form}, and $S^*AS\succ0$ by part (ii).
\item[$(vi)$] ($\Rightarrow$) If $A\oplus B\succ0$, then $(\forall\,\mathbf{x}\in\mathbb{F}^n\setminus\{\mathbf{0}\})\,q_A(\mathbf{x})=q_{A\oplus B}(\mathbf{x}\oplus\mathbf{0})>0$.\\
Likewise, $(\forall\,\mathbf{y}\in\mathbb{F}^m\setminus\{\mathbf{0}\})\,q_B(\mathbf{y})=q_{A\oplus B}(\mathbf{0}\oplus\mathbf{y})>0$, so $A\succ0$ and $B\succ0$.
\item[$(vi)$] ($\Leftarrow$) If $A\succ0$, $B\succ0$, and $(\mathbf{x},\mathbf{y})\neq(\mathbf{0},\mathbf{0})$, then $q_{A\oplus B}(\mathbf{x}\oplus\mathbf{y})=q_A(\mathbf{x})+q_B(\mathbf{y})>0$. Thus $A\oplus B\succ0$ by part (iii).
\end{enumerate}
\hlt{Structural Commentary}: The identities for $q_M$ preserve nonnegativity and expose the zero sets of each construction.\\
\hlt{Key Insight}: Strict positivity is the absence of nonzero vectors in the transformed zero set.
\end{proof}
\end{theorem}

\begin{proposition} \label{prop:two-by-two-positive-semidefinite-square-root} \hlt{Explicit Square Root for a $2\times2$ Positive Semidefinite Matrix}\\
Let $\mathbb{F}\in\{\mathbb{R},\mathbb{C}\}$ and $A\in\operatorname{Mat}(2,2,\mathbb{F})$ satisfy $A\succeq0$ and $A\neq0$. Set $\tau\defeq\bigl(\operatorname{tr}(A)+2\sqrt{\det(A)}\bigr)^{1/2}$.
\begin{enumerate}[label=(\roman*)]
\setlength{\itemsep}{0pt}
\item $\tau>0$.
\item $A^{1/2}=\frac{1}{\tau}\bigl(A+\sqrt{\det(A)}\,I_2\bigr)$.
\end{enumerate}
\begin{proof}
\begin{enumerate}[label=(\roman*)]
\setlength{\itemsep}{0pt}
\item $\det(A)\geq0$ and $\operatorname{tr}(A)>0$ by Proposition \ref{prop:trace-determinant-positive-semidefinite}(ii),(iii), so $\tau>0$.
\item Set $C\defeq\frac{1}{\tau}\bigl(A+\sqrt{\det(A)}\,I_2\bigr)$. Then $C\succeq0$ by Theorem \ref{thm:positive-semidefinite-stability-rules}(i), since $I_2\succeq0$ and $\tau^{-1},\tau^{-1}\sqrt{\det(A)}\geq0$.\\
The identity $A^2-\operatorname{tr}(A)A+\det(A)I_2=0$ holds by Theorems \ref{thm:matrix-characteristic-polynomial-basic-coefficients} and \ref{thm:cayley-hamilton-matrix}, regarding a real $A$ as complex when necessary.\\
Thus $\tau^2C^2=A^2+2\sqrt{\det(A)}A+\det(A)I_2=\tau^2A$, so $C^2=A$ and $C=A^{1/2}$ by Corollary \ref{cor:positive-semidefinite-square-root-formula}(i).
\end{enumerate}
\hlt{Structural Commentary}: Cayley--Hamilton reduces the squared linear candidate, and positive-root uniqueness identifies it.\\
\hlt{Key Insight}: $\tau^2=\operatorname{tr}(A)+2\sqrt{\det(A)}$ normalizes the square to $A$.
\end{proof}
\end{proposition}

\begin{definition} \label{def:directed-incidence-matrix} \hlt{Directed Incidence Matrix}\\
Let $m,n\in\mathbb{Z}_{>0}$, and let $G=(V,E)$ be a directed graph with $V=\{1,\ldots,n\}$ and ordered edges $E=(e_k)_{k=1}^m$, where $e_k=u_k\to v_k$.\\
Its \hlt{directed incidence matrix} is the matrix $B\in\operatorname{Mat}(m,n,\mathbb{R})$ satisfying $(\forall\,k\in\{1,\ldots,m\})\,B_{k,\cdot}=-\mathbf{e}_{u_k}^t+\mathbf{e}_{v_k}^t$.
\end{definition}

Figure \ref{fig:directed-incidence-matrix} records the chosen edge orientations and their corresponding signed incidence rows.

\begin{illustrationfigure}[H]
  \centering
  \begin{tikzpicture}[
  >=Stealth,
  every node/.style={font=\small},
  vertex/.style={circle, draw=blue!65!black, fill=blue!7, minimum size=0.62cm},
  edge/.style={->, blue!70!black, line width=0.75pt},
  mat/.style={draw=gray!55, fill=gray!5, rounded corners=2pt, inner sep=5pt, align=center}
]
  \node[vertex] (v1) at (0,1.2) {$1$};
  \node[vertex] (v2) at (2.0,1.2) {$2$};
  \node[vertex] (v3) at (2.0,-0.6) {$3$};
  \node[vertex] (v4) at (0,-0.6) {$4$};

  \draw[edge] (v1) -- node[above] {$e_1$} (v2);
  \draw[edge] (v2) -- node[right] {$e_2$} (v3);
  \draw[edge] (v4) -- node[below] {$e_3$} (v3);
  \draw[edge] (v1) -- node[left] {$e_4$} (v4);

  \node[mat] at (5.0,0.3) {$B=\begin{bmatrix}
  -1&1&0&0\\
  0&-1&1&0\\
  0&0&1&-1\\
  -1&0&0&1
  \end{bmatrix}$};
\end{tikzpicture}

  \caption{A directed incidence matrix records oriented edge differences from vertices to edges.}
  \label{fig:directed-incidence-matrix}
\end{illustrationfigure}

\begin{example} \label{ex:directed-incidence-matrix-four-node} \hlt{A Directed Incidence Matrix}\\
For Figure \ref{fig:directed-incidence-matrix}, the directed incidence matrix is
\begin{displaymath}
B=
\begin{bmatrix}
-1 & 1 & 0 & 0\\
0 & -1 & 1 & 0\\
0 & 0 & 1 & -1\\
-1 & 0 & 0 & 1
\end{bmatrix}.
\end{displaymath}
For node potentials $\mathbf{x}=(x_1,x_2,x_3,x_4)^t$, the edge-difference vector is
\begin{displaymath}
B\mathbf{x}=
\begin{bmatrix}
x_2-x_1\\
x_3-x_2\\
x_3-x_4\\
x_4-x_1
\end{bmatrix}.
\end{displaymath}
The equivalence $B\mathbf{x}=\mathbf{0}\Longleftrightarrow x_1=x_2=x_3=x_4$ implies $\ker(B)=\operatorname{span}\{(1,1,1,1)^t\}$.\\
Thus $\operatorname{rank}(B)=3$ by Corollary \ref{cor:matrix-rank-nullity}(i). For $\mathbf{y}\defeq(1,1,-1,-1)^t$, $B^t\mathbf{y}=\mathbf{0}$ by direct multiplication.\\
The equality $\operatorname{rank}(B^t)=3$ holds by Theorem \ref{thm:row-rank-equals-column-rank}, so $\ker(B^t)=\operatorname{span}\{\mathbf{y}\}$ by Corollary \ref{cor:matrix-rank-nullity}(i).\\
Thus $\mathbf{b}\in\operatorname{col}(B)$ iff $\mathbf{y}^t\mathbf{b}=b_1+b_2-b_3-b_4=0$ by Proposition \ref{prop:adjoint-kernel-range-orthogonality}(iv). This is the Kirchhoff loop-compatibility condition.\\
\hlt{Structural Commentary}: The kernel records constant-potential freedom; the left kernel records edge-data compatibility.\\
\hlt{Key Insight}: The oriented rows encode both edge differences and the cycle relation that determines which edge data are realizable.
\end{example}

\begin{perspective} \label{per:incidence-least-squares-ranking} \hlt{Incidence Data and Least Squares}\\
For an incidence matrix $B$ as in Example \ref{ex:directed-incidence-matrix-four-node}, $B\mathbf{x}=\mathbf{b}$ seeks node potentials reproducing prescribed edge differences.\\
If loop compatibility fails, a least-squares solution projects $\mathbf{b}$ onto $\operatorname{col}(B)$ and satisfies $B^*(B\mathbf{x}_0-\mathbf{b})=\mathbf{0}$ by Theorem \ref{thm:matrix-least-squares-normal-equations}.\\
The constant-vector kernel makes the potentials unique only after fixing a reference node or selecting the minimum-norm solution.
\end{perspective}

\begin{proposition} \label{prop:incidence-laplacian-positive-semidefinite} \hlt{Incidence Laplacians Are Positive Semidefinite}\\
Let $G=(V,E)$ be a directed graph, let $B\in\operatorname{Mat}(m,n,\mathbb{R})$ be its directed incidence matrix, and set $L\defeq B^tB$.\\
Let $G_{\mathrm{u}}$ be the underlying undirected graph obtained by forgetting edge orientations.
\begin{enumerate}[label=(\roman*)]
\setlength{\itemsep}{0pt}
\item $L\succeq0$.
\item $(\forall\,\mathbf{x}\in\mathbb{R}^n)\,\mathbf{x}^tL\mathbf{x}=\sum\limits_{e=(u\to v)\in E}(x_v-x_u)^2$.
\item $\ker(L)$ is the subspace of vectors that are constant on each connected component of $G_{\mathrm{u}}$.
\end{enumerate}
\begin{proof}
\begin{enumerate}[label=(\roman*)]
\setlength{\itemsep}{0pt}
\item $L=B^*B\succeq0$ by Proposition \ref{prop:rectangular-gram-products-are-psd}(i).
\item For $\mathbf{x}\in\mathbb{R}^n$, $(\forall\,e=(u\to v)\in E)\,(B\mathbf{x})_e=x_v-x_u$ by the incidence-row formula. Thus $\mathbf{x}^tL\mathbf{x}=\norm{B\mathbf{x}}_2^2=\sum\limits_{e=(u\to v)\in E}(x_v-x_u)^2$.
\item $\ker(L)=\ker(B)$ by Theorem \ref{thm:gram-factor-kernel-rank-column-data}(ii). For $\mathbf{x}\in\mathbb{R}^n$, $B\mathbf{x}=\mathbf{0}$ iff $x_u=x_v$ on every edge.\\
These equalities propagate along paths and are equivalent to componentwise constancy on $G_{\mathrm{u}}$.
\end{enumerate}
\hlt{Structural Commentary}: The Gram identity converts the Laplacian form and kernel into edge-difference conditions.\\
\hlt{Key Insight}: Zero edge differences are exactly componentwise-constant vertex values.
\end{proof}
\end{proposition}
